\chapter{UVM Verification Environment}
\label{ch:uvm-env}

% DRAFT BRIEF: walk the testbench bottom-up. Write structure now; fill 7.8 results after
% regressions are green. Format = component diagrams + waveform screenshots; avoid excessive
% SV listings. RE-ENUMERATE vseqs from disk at write time; footnote that the 56-vseq suite is
% a WIP snapshot (~105-case testplan target).
% Sources: tb_i3c_top.sv, i3c_scoreboard.sv, i3c_driver.sv, i3c_monitor.sv, i3c_base_vseq.sv,
% i3c_vseq_list.sv, Makefile, i3c_vseqs/**, sva/**, verification specs 06-09.

\section{TB top and interfaces}
\brief{\texttt{tb_i3c_top} instantiates the DUT + reg/i3c interfaces, clock/reset, the
\SI{100}{\milli\second} watchdog \texttt{uvm_fatal}, SHM dump gated by \texttt{+DUMP_WAVES},
and the SVA binds. Note reduced TB FIFO depths vs RTL to shrink sim time.}

\section{Register agent}
\brief{\texttt{dv_reg/} agent driving the 32-bit addr/wen/ren bus
(driver/monitor/sequencer/cfg).}

\figph{Register agent}{reg agent driver/monitor/sequencer on the 32-bit bus}{dv_reg/}{TikZ}

\section{I3C device-mode agent (driver + monitor)}
\brief{\texttt{dv_i3c/} device-mode responder: \texttt{i3c_driver} with 13 device-only phases
(\texttt{i3c_drv_phase_e}), CCC-aware \texttt{i3c_monitor} capturing protocol events. Describe
the phase FSM and the monitor's protocol state machine.}

\figph{I3C agent (device-mode driver + monitor)}{driver 13-phase FSM + monitor protocol capture}
  {dv_i3c/i3c_agent_pkg.sv, i3c_driver.sv, i3c_monitor.sv}{TikZ}
\tabph{Driver 13-phase table (DrvIdle...DrvDAA)}{i3c_agent_pkg.sv}

\section{Env / vseqr / scoreboard class hierarchy}
\brief{\texttt{i3c_env} (reg agent + i3c agent + scoreboard + virtual sequencer), the class
hierarchy, and \texttt{uvm_config_db} wiring. Reference the scoreboard checks from
\cref{ch:verif-method}.}

\figph{Env class hierarchy}
  {i3c_base_test to i3c_env to \{reg_agent, i3c_agent, scoreboard, virtual_sequencer\}}
  {redraw from print_topology}{TikZ}
\tabph{Scoreboard check list}{i3c_scoreboard.sv}

\section{Vseq library}
\brief{Re-enumerate from disk. Snapshot at planning time: 56 vseqs --- csr 14 / fifo 5 /
bus 12 / sdrw 9 / sdrr 10 / resp 4 / ccc 1 / imm 1. CCC vseq is ad-hoc only
(\texttt{make sim SEQ=i3c_ccc_broadcast_enec_vseq}). Footnote the WIP snapshot vs
$\sim$105-case testplan.}

\tabph{Vseq library by category (re-enumerate count at write time)}{i3c_vseqs/**}

\section{Build/run flow}
\brief{Makefile targets: \texttt{compile}, \texttt{sim SEQ=...}, the eight category regression
targets (smoke, regression, csr/fifo/bus/sdrw/sdrr/sdr regressions), \texttt{COV=1} coverage,
\texttt{DUMP_WAVES=1}, and the \texttt{fifo_non_power_of_two_elaboration} (FIFO_005)
elaboration-fail-as-pass guard. Note the env must be sourced (\texttt{XCELIUM1803.sh}) first.}

\tabph{Regression targets + SEQ lists}{Makefile}

\section{Waveform and debug (SimVision)}
\brief{SHM dump workflow, opening \texttt{waves/<SEQ>/waves.shm} in SimVision, the events
worth annotating (START/ACK/T-bit/OD-to-PP/STOP).}

\section{Regression results}
\brief{BLOCKED on green regressions. Insert the pass/fail matrix + representative waveforms
once \texttt{make regression} + category regressions pass.}

\figph{Representative write waveform}{SDR private write, annotated START/ACK/T-bit/STOP}
  {SimVision capture --- AFTER SIMS}{SimVision PDF/PNG}
\figph{Representative read waveform}{SDR private read, annotated T-bit end-of-data}
  {SimVision capture --- AFTER SIMS}{SimVision PDF/PNG}

\section{Phase-2 roadmap}
\brief{Known gaps + next steps: fix monitor stub if present, ENTDAA/CCC test breadth, error
injection, multi-device, functional coverage. Cross-reference \cref{ch:conclusion}.}

\figph{SVA binding map}{5 SVA modules to bound vs instantiated targets in tb_i3c_top}
  {sva/**, tb_i3c_top.sv}{TikZ}
\tabph{Five SVA modules (bind vs instantiate)}{sva/**, tb_i3c_top.sv}
